% Generated draft: requires mathematical and source review.
\documentclass[11pt]{article}
\usepackage[margin=25mm]{geometry}
\usepackage{fontspec}
\setmainfont{Latin Modern Roman}
\usepackage{amsmath,amssymb,mathtools,unicode-math}
\setmathfont{Latin Modern Math}
\usepackage{enumitem,needspace,xurl,hyperref}
\setlength{\parindent}{0pt}
\setlength{\parskip}{5pt}
\setlength{\emergencystretch}{3em}
\newcommand{\N}{\mathbb{N}}
\newcommand{\Z}{\mathbb{Z}}
\newcommand{\Q}{\mathbb{Q}}
\newcommand{\R}{\mathbb{R}}
\newcommand{\C}{\mathbb{C}}
\newcommand{\F}{\mathbb{F}}
\newcommand{\A}{\mathbb{A}}
\newcommand{\PP}{\mathbb{P}}
\newcommand{\E}{\mathbb{E}}
\newcommand{\eps}{\varepsilon}
\newcommand{\dd}{\,\mathrm{d}}
\begin{document}
\section*{UnsolvedMath ID: 30003962}
\textbf{Source problem code:} catalogue entry
\subsection{Definitions and mathematical statement}
The supplied problem concerns strata of holomorphic differentials, but does not specify the genus, zero multiplicities, base field, connected component, or whether differentials are projectivized. The following gives standard notation over \(\C\) to describe the questions; this choice of setting requires confirmation against the source.

For an integer \(g\geq 1\), let \(\mu=(m_1,\ldots,m_n)\) be a list of positive integers satisfying \(\sum_{i=1}^n m_i=2g-2\), allowing the empty list when \(g=1\). Write \(\mathcal H(\mu)\) for the moduli space of tuples \((X,\omega,p_1,\ldots,p_n)\), up to isomorphism, where \(X\) is a smooth connected projective complex curve of genus \(g\), \(\omega\) is a nonzero holomorphic differential, the \(p_i\) are distinct, and
\[
\operatorname{div}(\omega)=\sum_{i=1}^n m_i p_i.
\]
This notation labels the zeros; the supplied record does not specify whether zeros are labeled or whether a moduli stack or a coarse moduli space is intended. Let
\[
\mathbb P\mathcal H(\mu)=\mathcal H(\mu)/\C^{\times},
\qquad c\cdot(X,\omega,p_1,\ldots,p_n)=(X,c\omega,p_1,\ldots,p_n).
\]
On the projectivized stratum, the tautological line bundle \(L\) has fiber \(\C\omega\) at \((X,[\omega],p_1,\ldots,p_n)\). Use the convention \(\eta=c_1(L)\); the source does not specify its sign convention or the intersection or cohomology theory in which this class is considered.

\begin{enumerate}
\item Determine whether a stratum contains a complete curve, meaning a closed integral algebraic subvariety of dimension one that is proper over the base field. The supplied wording leaves unresolved whether the ambient space is \(\mathcal H(\mu)\), \(\mathbb P\mathcal H(\mu)\), or a chosen connected component. It also does not specify whether the question asks for existence in some stratum or a classification across strata.
\item Determine the positive integers \(k\) for which
\[
\eta^k=0
\]
on the intended stratum. A precise vanishing question requires specification of the projectivization, the ambient component, the ring containing \(\eta^k\), and its coefficient system. If an unprojectivized stratum is intended, the meaning of its tautological class must also be specified.
\item For an intended compactification \(\overline S\) of a stratum \(S\), determine positivity properties of the other divisor classes under consideration. The compactification, divisor classes, and meaning of positivity are not supplied; possible distinct notions include ampleness, nefness, bigness, and pseudoeffectivity. Thus this part remains a research question requiring these choices, rather than a specified proposition.
\end{enumerate}

The supplied dated curation reports that complete curves are ruled out in unprojectivized holomorphic strata, while the projectivized complete-curve question and the requested vanishing and positivity questions remain open. This is a supplied literature assessment, not an independently verified theorem or status determination.
\subsection{Short English statement}
The problem asks about complete curves in strata of holomorphic differentials, vanishing powers of their tautological class \(\eta\), and positivity of divisor classes on compactifications. Its precise scope requires clarification of projectivization, strata and components, the ring used for vanishing, and the compactification and positivity notion. The supplied open-status assessment is dated and requires independent review.
\subsection{Sources}
\begin{itemize}
\item Supplied local UnsolvedMath registry record, permanent ID \(30003962\), title: \emph{Complete Curves and Positivity in Differential Strata}; exact source problem code: \texttt{catalogue entry}. Collection: Oberwolfach Reports Open Problems. The record gives proposed year 2018 and supplies no named proposer.
\item Supplied source citation: \emph{Flat Surfaces and Algebraic Curves} (2019), Oberwolfach Reports, pp.\ 2583--2650, DOI \texttt{10.4171/owr/2018/42}. \url{https://doi.org/10.4171/owr/2018/42}. The supplied background describes Oberwolfach Reports as records of mathematical workshops and their open-problem discussions.
\item Supplied literature reference: ``Oberwolfach Report 42, strata problems.'' \url{https://doi.org/10.4171/owr/2018/42}.
\item Supplied literature reference: ``Chen, complete curves in strata of differentials.'' \url{https://doi.org/10.5802/crmath.416}. No fuller bibliographic citation or theorem statement is supplied.
\item Literature assessment checked on \texttt{2026-08-21}, according to the supplied record. This is supplied dated curation, not independently verified: it labels the problem open and reports nonexistence of complete curves in unprojectivized holomorphic strata, while leaving the projectivized complete-curve question and the vanishing and positivity questions open. The background explicitly requests independent verification of time-sensitive status and cited claims.
\item The registry labels the statement \texttt{corrected\_verified} and reports a manual semantic, typographic, and LaTeX review. These are supplied editorial assertions, not independent mathematical verification of this draft.
\end{itemize}
\end{document}
